
\documentclass{article}
\usepackage{colortbl}
\usepackage{makecell}
\usepackage{multirow}
\usepackage{supertabular}

\begin{document}

\newcounter{utterance}

\centering \large Interaction Transcript for game `dond', experiment `coop', episode 3 with claude{-}opus{-}4{-}8{-}azure.
\vspace{24pt}

{ \footnotesize  \setcounter{utterance}{1}
\setlength{\tabcolsep}{0pt}
\begin{supertabular}{c@{$\;$}|p{.15\linewidth}@{}p{.15\linewidth}p{.15\linewidth}p{.15\linewidth}p{.15\linewidth}p{.15\linewidth}}
    \# & $\;$A & \multicolumn{4}{c}{Game Master} & $\;\:$B\\
    \hline

    \theutterance \stepcounter{utterance}  
    & & \multicolumn{4}{p{0.6\linewidth}}{
        \cellcolor[rgb]{0.9,0.9,0.9}{
            \makecell[{{p{\linewidth}}}]{
                \texttt{\tiny{[P1$\langle$GM]}}
                \texttt{Du spelar ett kooperativt förhandlingsspel där du måste komma överens med en annan spelare om hur du delar upp en uppsättning föremål mellan er.} \\
\\ 
\texttt{Regler:} \\
\texttt{(a) Du och den andra spelaren får en uppsättning föremål. Du får också en hemlig värderfunktion som representerar hur mycket du värderar varje objekttyp.} \\
\texttt{(b) Du utbyter meddelanden med den andra spelaren för att komma överens om vem som får vilka föremål. Du kan skicka högst 5 meddelanden vardera eller avsluta spelet tidigt genom att göra ett hemligt förslag när som helst.} \\
\texttt{(c) Du ombeds att lämna in ett hemligt förslag som anger de föremål du vill ha, formaterat inom hakparenteser enligt följande: "{[}Förslag: <nummer> <objektnamn>, <nummer> <objektnamn>, <...>{]}"} \\
\texttt{(d) Om dina förslag är komplementära, det vill säga det finns tillräckligt med föremål för att uppfylla båda förslagen, får varje spelare en poäng baserad på summan av värdena för de föremål han fick. Annars får ni båda noll poäng.} \\
\texttt{(e) Ditt mål är att maximera summan av din poäng och den andra spelarens poäng.} \\
\\ 
\texttt{Låt oss börja.} \\
\\ 
\texttt{Uppsättningen av tillgängliga föremål är:} \\
\\ 
\texttt{1 tång, 2 tegelstenar, 1 ring, 1 klocka} \\
\\ 
\texttt{Din hemliga värderfunktion är:} \\
\\ 
\texttt{tång: 1, tegelsten: 0, ring: 2, klocka: 7} \\
\\ 
\texttt{VIKTIGT: Dina meddelanden, om det inte är ett hemligt förslag, överfördes direkt till den andra spelaren, så inkludera inget svar på reglerna eller text som meddelar ditt meddelande. Använd det angivna formatet för att göra ett hemligt förslag. Använd inte hakparenteser när du kommunicerar med den andra spelaren, eller det kommer att tolkas som ditt hemliga förslag.} \\
            }
        }
    }
    & \\ \\

    \theutterance \stepcounter{utterance}  
    & \multicolumn{4}{p{0.6\linewidth}}{
        \cellcolor[rgb]{0.9,0.9,0.9}{
            \makecell[{{p{\linewidth}}}]{
                \texttt{\tiny{[P1$\rangle$GM]}}
                \texttt{Hej! Låt oss försöka dela upp föremålen så att vi båda blir nöjda. För min del värderar jag klockan högst, medan tegelstenarna inte betyder så mycket för mig. Vilka föremål är viktigast för dig? Om du berättar kan vi hitta en fördelning som maximerar vårt gemensamma värde.} \\
            }
        }
    }
    & & \\ \\

    \theutterance \stepcounter{utterance}  
    & & & \multicolumn{2}{p{0.3\linewidth}}{
        \cellcolor[rgb]{0.9,0.9,0.9}{
            \makecell[{{p{\linewidth}}}]{
                \texttt{\tiny{[GM$|$GM]}}
                \texttt{continue} \\
            }
        }
    }
    & & \\ \\

    \theutterance \stepcounter{utterance}  
    & & & \multicolumn{2}{p{0.3\linewidth}}{
        \cellcolor[rgb]{0.9,0.9,0.9}{
            \makecell[{{p{\linewidth}}}]{
                \texttt{\tiny{[GM$|$GM]}}
                \texttt{Hej! Låt oss försöka dela upp föremålen så att vi båda blir nöjda. För min del värderar jag klockan högst, medan tegelstenarna inte betyder så mycket för mig. Vilka föremål är viktigast för dig? Om du berättar kan vi hitta en fördelning som maximerar vårt gemensamma värde.} \\
            }
        }
    }
    & & \\ \\

    \theutterance \stepcounter{utterance}  
    & & \multicolumn{4}{p{0.6\linewidth}}{
        \cellcolor[rgb]{0.9,0.9,0.9}{
            \makecell[{{p{\linewidth}}}]{
                \texttt{\tiny{[GM$\rangle$P2]}}
                \texttt{Du spelar ett kooperativt förhandlingsspel där du måste komma överens med en annan spelare om hur du delar upp en uppsättning föremål mellan er.} \\
\\ 
\texttt{Regler:} \\
\texttt{(a) Du och den andra spelaren får en uppsättning föremål. Du får också en hemlig värderfunktion som representerar hur mycket du värderar varje objekttyp.} \\
\texttt{(b) Du utbyter meddelanden med den andra spelaren för att komma överens om vem som får vilka föremål. Du kan skicka högst 5 meddelanden vardera eller avsluta spelet tidigt genom att göra ett hemligt förslag när som helst.} \\
\texttt{(c) Du ombeds att lämna in ett hemligt förslag som anger de föremål du vill ha, formaterat inom hakparenteser enligt följande: "{[}Förslag: <nummer> <objektnamn>, <nummer> <objektnamn>, <...>{]}"} \\
\texttt{(d) Om dina förslag är komplementära, det vill säga det finns tillräckligt med föremål för att uppfylla båda förslagen, får varje spelare en poäng baserad på summan av värdena för de föremål han fick. Annars får ni båda noll poäng.} \\
\texttt{(e) Ditt mål är att maximera summan av din poäng och den andra spelarens poäng.} \\
\\ 
\texttt{Låt oss börja.} \\
\\ 
\texttt{Uppsättningen av tillgängliga föremål är:} \\
\\ 
\texttt{1 tång, 2 tegelstenar, 1 ring, 1 klocka} \\
\\ 
\texttt{Din hemliga värderfunktion är:} \\
\\ 
\texttt{tång: 1, tegelsten: 1, ring: 4, klocka: 3} \\
\\ 
\texttt{VIKTIGT: Dina meddelanden, om det inte är ett hemligt förslag, överfördes direkt till den andra spelaren, så inkludera inget svar på reglerna eller text som meddelar ditt meddelande. Använd det angivna formatet för att göra ett hemligt förslag. Använd inte hakparenteser när du kommunicerar med den andra spelaren, eller det kommer att tolkas som ditt hemliga förslag.} \\
\\ 
\texttt{Hej! Låt oss försöka dela upp föremålen så att vi båda blir nöjda. För min del värderar jag klockan högst, medan tegelstenarna inte betyder så mycket för mig. Vilka föremål är viktigast för dig? Om du berättar kan vi hitta en fördelning som maximerar vårt gemensamma värde.} \\
            }
        }
    }
    & \\ \\

    \theutterance \stepcounter{utterance}  
    & & & \multicolumn{4}{p{0.6\linewidth}}{
        \cellcolor[rgb]{0.9,0.9,0.9}{
            \makecell[{{p{\linewidth}}}]{
                \texttt{\tiny{[GM$\langle$P2]}}
                \texttt{Tack för informationen! För mig är ringen det klart viktigaste föremålet, följt av klockan. Eftersom du värderar klockan högst känns det logiskt att du får klockan och jag får ringen.} \\
\\ 
\texttt{Det som återstår är tången och de två tegelstenarna, som jag värderar ganska lågt (1 vardera). Hur ser dina värden ut för dem? Om du värderar dem högre än mig bör du ta dem, så vi maximerar vårt gemensamma värde.} \\
            }
        }
    }
     \\ \\

    \theutterance \stepcounter{utterance}  
    & & & \multicolumn{2}{p{0.3\linewidth}}{
        \cellcolor[rgb]{0.9,0.9,0.9}{
            \makecell[{{p{\linewidth}}}]{
                \texttt{\tiny{[GM$|$GM]}}
                \texttt{continue} \\
            }
        }
    }
    & & \\ \\

    \theutterance \stepcounter{utterance}  
    & & & \multicolumn{2}{p{0.3\linewidth}}{
        \cellcolor[rgb]{0.9,0.9,0.9}{
            \makecell[{{p{\linewidth}}}]{
                \texttt{\tiny{[GM$|$GM]}}
                \texttt{Tack för informationen! För mig är ringen det klart viktigaste föremålet, följt av klockan. Eftersom du värderar klockan högst känns det logiskt att du får klockan och jag får ringen.} \\
\\ 
\texttt{Det som återstår är tången och de två tegelstenarna, som jag värderar ganska lågt (1 vardera). Hur ser dina värden ut för dem? Om du värderar dem högre än mig bör du ta dem, så vi maximerar vårt gemensamma värde.} \\
            }
        }
    }
    & & \\ \\

    \theutterance \stepcounter{utterance}  
    & & \multicolumn{4}{p{0.6\linewidth}}{
        \cellcolor[rgb]{0.9,0.9,0.9}{
            \makecell[{{p{\linewidth}}}]{
                \texttt{\tiny{[P1$\langle$GM]}}
                \texttt{Tack för informationen! För mig är ringen det klart viktigaste föremålet, följt av klockan. Eftersom du värderar klockan högst känns det logiskt att du får klockan och jag får ringen.} \\
\\ 
\texttt{Det som återstår är tången och de två tegelstenarna, som jag värderar ganska lågt (1 vardera). Hur ser dina värden ut för dem? Om du värderar dem högre än mig bör du ta dem, så vi maximerar vårt gemensamma värde.} \\
            }
        }
    }
    & \\ \\

    \theutterance \stepcounter{utterance}  
    & \multicolumn{4}{p{0.6\linewidth}}{
        \cellcolor[rgb]{0.9,0.9,0.9}{
            \makecell[{{p{\linewidth}}}]{
                \texttt{\tiny{[P1$\rangle$GM]}}
                \texttt{Det låter som en bra fördelning! Du får ringen och jag får klockan.} \\
\\ 
\texttt{För tången är vi lika (jag värderar den till 1), så den kan jag ta. De två tegelstenarna värderar du högre än jag (jag har 0 på dem), så du bör ta dem.} \\
\\ 
\texttt{Så förslaget blir: du får ringen och de två tegelstenarna, jag får klockan och tången. Låter det bra?} \\
\\ 
\texttt{{[}Förslag: 1 tång, 1 klocka{]}} \\
            }
        }
    }
    & & \\ \\

    \theutterance \stepcounter{utterance}  
    & & & \multicolumn{2}{p{0.3\linewidth}}{
        \cellcolor[rgb]{0.9,0.9,0.9}{
            \makecell[{{p{\linewidth}}}]{
                \texttt{\tiny{[GM$|$GM]}}
                \texttt{proposal} \\
            }
        }
    }
    & & \\ \\

    \theutterance \stepcounter{utterance}  
    & & & \multicolumn{2}{p{0.3\linewidth}}{
        \cellcolor[rgb]{0.9,0.9,0.9}{
            \makecell[{{p{\linewidth}}}]{
                \texttt{\tiny{[GM$|$GM]}}
                [1, 0, 0, 1]
            }
        }
    }
    & & \\ \\

    \theutterance \stepcounter{utterance}  
    & & \multicolumn{4}{p{0.6\linewidth}}{
        \cellcolor[rgb]{0.9,0.9,0.9}{
            \makecell[{{p{\linewidth}}}]{
                \texttt{\tiny{[GM$\rangle$P2]}}
                \texttt{Den andra spelaren har lämnat in ett hemligt förslag. Det är nu din tur att lämna in ett eget förslag.} \\
            }
        }
    }
    & \\ \\

    \theutterance \stepcounter{utterance}  
    & & & \multicolumn{4}{p{0.6\linewidth}}{
        \cellcolor[rgb]{0.9,0.9,0.9}{
            \makecell[{{p{\linewidth}}}]{
                \texttt{\tiny{[GM$\langle$P2]}}
                \texttt{{[}Förslag: 1 ring{]}} \\
            }
        }
    }
     \\ \\

    \theutterance \stepcounter{utterance}  
    & & & \multicolumn{2}{p{0.3\linewidth}}{
        \cellcolor[rgb]{0.9,0.9,0.9}{
            \makecell[{{p{\linewidth}}}]{
                \texttt{\tiny{[GM$|$GM]}}
                \texttt{proposal} \\
            }
        }
    }
    & & \\ \\

    \theutterance \stepcounter{utterance}  
    & & & \multicolumn{2}{p{0.3\linewidth}}{
        \cellcolor[rgb]{0.9,0.9,0.9}{
            \makecell[{{p{\linewidth}}}]{
                \texttt{\tiny{[GM$|$GM]}}
                [0, 0, 1, 0]
            }
        }
    }
    & & \\ \\

    \theutterance \stepcounter{utterance}  
    & & & \multicolumn{2}{p{0.3\linewidth}}{
        \cellcolor[rgb]{0.9,0.9,0.9}{
            \makecell[{{p{\linewidth}}}]{
                \texttt{\tiny{[GM$|$GM]}}
                [[1, 0, 0, 1], [0, 0, 1, 0]]
            }
        }
    }
    & & \\ \\

\end{supertabular}
}

\end{document}
